% ============================================================
% 50_typography_semantics.tex — Schriftmakros für Box-Typen & semantische Inhaltsmakros
% ============================================================

% Explizite bp-Grössen (wie Analysis) für präzise, font-unabhängige Kontrolle
\newcommand{\ZSFfontChapter}{\sffamily\bfseries\fontsize{8.1bp}{9.1bp}\selectfont}
\newcommand{\ZSFfontSection}{\sffamily\bfseries\fontsize{8.1bp}{9.1bp}\selectfont}
\newcommand{\ZSFfontSubsubsection}{\sffamily\bfseries\fontsize{7.2bp}{8.0bp}\selectfont}
\newcommand{\ZSFfontBoxTitle}{\sffamily\bfseries\fontsize{8.1bp}{9.1bp}\selectfont}
% chktex 6
\newcommand{\ZSFfontNote}{\sffamily\fontsize{6.7bp}{7.5bp}\itshape\selectfont}
% Rechtsbuendiges Meta-/Kategorie-Tag im Titelbalken (vgl. \ZSFTitleTag in
% styles/60_boxes.tex). Etwas kleiner und kursiv, damit es den Titel begleitet
% statt mit ihm zu konkurrieren.
\newcommand{\ZSFfontTitleTag}{\sffamily\fontsize{6.4bp}{7.0bp}\itshape\selectfont}
\newcommand{\ZSFfontTitleHeader}{\sffamily\fontsize{11bp}{12.5bp}\bfseries\selectfont}
\newcommand{\ZSFfontFooter}{\sffamily\bfseries\fontsize{9.5pt}{9.5pt}\selectfont}
\newcommand{\ZSFfontIdentityMarker}{\fontsize{4pt}{4pt}\selectfont}

% Schrift/Farbe für Code-Kommentare — nur für das optionale Code-Modul
% (styles/67_code_comments.tex) benötigt. ce_green stammt aus
% styles/40_colors_structure.tex. Schadlos wenn das Modul deaktiviert ist.
\newcommand{\ZSFfontCodeComment}{\ttfamily\color{ce_green}}

% Semantische Inhalts-Fontgrössen — die zwei Stufen des Reglers `font`
% (`normal` | `dense`). Kapitel wählen sie nicht direkt, sondern über den
% Regler an der Box (styles/60_boxes.tex) oder an der ZSFtable
% (styles/20_tables.tex). `dense` ist die eine Stufe kleiner: lange Register,
% Vergleichsmatrizen, gedrängte Aufzählungen.
%
% EIN Paar für beide Bausteine, nicht je eines: »eine Stufe kleiner« ist
% dieselbe Entscheidung, ob sie eine Tabelle oder einen Boxinhalt trifft. Mit
% zwei Paaren müsste jede Änderung am einen im anderen nachgezogen werden
% (styles/README → Vereinigungstest).
%
% Beide holen die Display-Skips zurück: \normalsize und \footnotesize setzen
% sie auf die Klassenwerte, und die ZSF-Entscheidung dazu ist eine andere
% (\ZSFApplyDisplaySkips in 30_layout_spacing). Ohne diesen Nachsatz wären
% Formeln in jeder Box luftiger als im Fliesstext — lautlos, weil beides
% für sich korrekt aussieht.
\newcommand{\ZSFfontContent}{\normalsize\ZSFApplyDisplaySkips}
\newcommand{\ZSFfontContentDense}{\footnotesize\ZSFApplyDisplaySkips}

%
%
%
%
%
%
%
%
%
\newcommand{\ZSFfontTableHead}{\bfseries}
% Marke einer ZSFlist und Beschriftung eines \ZSFsep-Blockwechsels — beide
% sind dieselbe Rolle: eine kurze fette Beschriftung innerhalb eines Bausteins.
\newcommand{\ZSFfontBlockLabel}{\bfseries}
% Text, der per \textVorBox/\textNachBox an eine Box gebunden wird.
\newcommand{\ZSFfontBoxBinding}{\small}
% Inline-Pille für Stolperfallen (\ZSFdanger).
\newcommand{\ZSFfontDangerTag}{\sffamily\bfseries}
%
%
%
\newcommand{\ZSFfontRefLocator}{\upshape\bfseries}

%
%
%
%
%
\newcommand{\ZSFBoxNote}[1]{{\ZSFfontNote\color{MutedTextColor}#1}}

%
%
%
%
%
%
%
%
%
%
%
%
%
%
%
%
%
%
%
%
%
%
%
%
%
%
\providecommand{\ZSFDiagramLabelScale}{1}
% Die Grösse wird gerechnet statt geschrieben. Über \the und damit in pt, weil
% \fontsize sein Argument sonst als Zahl ohne Einheit lesen müsste — der Faktor
% kommt so als fertiges Mass an und nicht als Term, den \fontsize zerlegt.
\makeatletter
\newlength{\ZSF@fontSize}
\newlength{\ZSF@fontSkip}
\newcommand{\ZSF@scaledFont}[3]{%
  \setlength{\ZSF@fontSize}{#1\dimexpr #2\relax}%
  \setlength{\ZSF@fontSkip}{#1\dimexpr #3\relax}%
  \fontsize{\the\ZSF@fontSize}{\the\ZSF@fontSkip}\selectfont}
\newcommand{\ZSFfontDiagramLabel}{%
  \sffamily\ZSF@scaledFont{\ZSFDiagramLabelScale}{5.8bp}{6.4bp}}
\newcommand{\ZSFfontDiagramLabelSmall}{%
  \sffamily\ZSF@scaledFont{\ZSFDiagramLabelScale}{4.6bp}{5.1bp}}
\makeatother

%
%
%
%
%
%
%
%
%
%
%
%
%
%
%
%
%
%
%
%
%
%
%
%
\NewDocumentCommand{\ZSFkeyword}{s o m}{\ZSFkeywordStyle{#3}}

%
%
%
%
%
\newcommand{\ZSFlabel}[1]{\ZSFlabelStyle{#1}}

% Die beiden Darstellungs-Hooks der Marker — getrennt, weil \ZSFkeyword und
% \ZSFlabel verschiedene Dinge bedeuten und eine ZSF den Unterschied ändern
% können soll, ohne ein einziges Kapitel anzufassen.
%
% BEIDE schlicht fett, ohne Farbe. \ZSFkeyword stand bis hierher fett in der
% KAPITELFARBE, mit dem Argument, der primäre Scan-Anker müsse sich von jeder
% anderen Fettung abheben. Das Argument ist am Satz gescheitert: In laufendem
% Text bringt eine eingefärbte Vokabel Unruhe in den Absatz, ohne dass die
% Farbe etwas sagt — sie wiederholt nur, in welchem Kapitel man ohnehin gerade
% ist. Wer den Absatz liest, wird gestört; wer ihn überfliegt, findet die
% Fettung auch ohne Farbe.
%
% Farbe bleibt damit im System das, was sie vorher schon war und wofür sie
% trägt: Kapitel-Identität auf Flächen (Balken, Akzente), Wegweiser zum ZIEL
% eines Verweises (\ZSFref) und semantische Zuordnung in Formeln (\ZSFmhl*).
% Als blosse Betonung im Fliesstext ist sie ausdrücklich nicht vorgesehen.
%
% Die Unterscheidung der beiden Marker liegt jetzt allein in der Bedeutung —
% \ZSFkeyword ist ein Fachbegriff und landet im Register, \ZSFlabel ist eine
% neutrale Beschriftung und tut das nicht. Das ist der Unterschied, der beim
% Schreiben zählt; er war nie der Farbe zu entnehmen.
\newcommand{\ZSFkeywordStyle}[1]{\textbf{#1}}
\newcommand{\ZSFlabelStyle}[1]{\textbf{#1}}

\newcommand{\ZSFScriptRef}[1]{\hfill\textcolor{\ScriptRefColor}{\ZSFfontBoxBinding (S.#1)}}

% \ZSFhl{...} — Hervorhebung des EINEN prüfungskritischsten Fakts pro Box
% (fett + unterstrichen). Bewusst sparsam: hoechstens einmal pro Block.
% Klare Rollentrennung:
%   * Fachbegriff/Scan-Anker        -> \ZSFkeyword
%   * Stolperfalle/kritische Ausnahme -> \ZSFdanger
%   * Folgerung                      -> \ZSFconclusion
%   * sonst die EINE Schlüsselaussage -> \ZSFhl
% \uline statt \underline: \underline steckt seinen Inhalt in eine
% unzerbrechliche Box — jede Hervorhebung, die länger als eine Zeile ist,
% sprengt damit die 50-mm-Spalte. \ZSFhl markiert aber laut Zweck eine ganze
% Aussage, nicht ein Wort. \uline bricht mit dem Absatz um.
% Fett als Deklaration, nicht als Argument von \uline: ein \textbf{...} als
% Argument macht den Inhalt für ulem wieder zu einer Einheit.
%
% Die Darstellung liegt wie bei \ZSFkeyword und \ZSFlabel in einem eigenen
% Hook. Ohne ihn erwischte »die Hervorhebung dieser ZSF umstellen« zwei von
% drei Markern — der dritte war der einzige mit fest verdrahteter Erscheinung,
% und man sah es erst am Satz (styles/README → Auszeichnungs-Hooks).
\newcommand{\ZSFhlStyle}[1]{{\bfseries\uline{#1}}}
\newcommand{\ZSFhl}[1]{\ZSFhlStyle{#1}}

% Papierorientierte Referenzrollen:
%   * \ZSFref{label}: hilfreicher Querverweis mit farbiger Abschnittsnummer.
%   * \ZSFsectionref{label}: kompakte Zielnummer für lokale Router/Tabellen.
% Ziel-Label via optionales Argument von \SubsubsectionBar/\SubsectionBar/\StartChapter setzen:
%   \SubsectionBar[sec:matrix-inverse]{Inverse}
% Farbe = Palette-Farbe des ZIELkapitels (Wegweiser-Prinzip): wird automatisch
% aus der Kapitelnummer der Referenz abgeleitet ("6.6" -> ChapterColor6).
% Bei Kapiteln jenseits der Palette wird derselbe Modulo-Lookup wie bei
% \StartChapter genutzt, damit gewrappte Kapitel und Verweise konsistent sind.
% Aufrecht + fett statt kursiv für maximale Lesbarkeit bei 6.7bp.
\makeatletter
\def\ZSF@refExtractFirst#1#2\@nil{#1}
\def\ZSF@refExtractChap#1.#2\@nil{#1}
% Bestimmt \ZSF@refTargetColor aus dem Label; Fallback auf die aktive
% Kapitelfarbe im ersten latexmk-Pass (Label noch undefiniert).
\newcommand{\ZSF@refComputeTargetColor}[1]{%
  \@ifundefined{r@#1}{%
    \edef\ZSF@refTargetColor{\chaptercolor}%
  }{%
    \edef\ZSF@refTargetNum{\expandafter\expandafter\expandafter
      \ZSF@refExtractFirst\csname r@#1\endcsname\@nil}%
    \edef\ZSF@refTargetChap{\expandafter\ZSF@refExtractChap\ZSF@refTargetNum.\@nil}%
    \ZSFSetPaletteColorNameFromNumber{\ZSF@refTargetChap}{\ZSF@refTargetColor}%
  }%
}
% Die Zielfarbe läuft über \ZSFInk. Ohne das steht ein Verweis auf das eigene
% Kapitel im Titel einer \SubsectionBar als ChapterColorN auf \chaptercolor!75 —
% derselbe Farbton bei voller Sättigung, also unlesbar, ohne dass Build oder
% Linter etwas melden. Auf einer Fläche mit eigener Kontrastfarbe erbt der
% Verweis diese und bleibt lesbar; der Farb-Wegweiser gilt dort, wo er auch
% ankommt — im Boxinhalt (40_colors_structure → Ink-Vertrag).
\newcommand{\ZSFref}[1]{%
  \begingroup
  \ZSF@refComputeTargetColor{#1}%
  \hyperref[#1]{{\ZSFfontNote\ZSFfontRefLocator
    \ZSFInk{\ZSF@refTargetColor}{($\rightarrow$\,\ref*{#1})}}}%
  \endgroup
}
\newcommand{\ZSFsectionref}[1]{%
  \begingroup
  \ZSF@refComputeTargetColor{#1}%
  \hyperref[#1]{{\ZSFfontRefLocator\ZSFInk{\ZSF@refTargetColor}{\ref*{#1}}}}%
  \endgroup
}
% \ZSFhubref{label} — Verweis mit Zielnummer UND Zielname. [FORK-ERWEITERUNG]
%
% Dritte Rolle derselben Familie, nicht eine dritte Schreibweise derselben
% Sache. Der Unterschied ist, wer den Kontext trägt:
%   \ZSFref        Nummer neben einer Aussage, die selbst schon sagt, worum es
%                  geht — der Verweis ist Zusatz.
%   \ZSFsectionref blosse Zielnummer in einer Tabellenzelle, deren Zeile den
%                  Kontext trägt.
%   \ZSFhubref     der Verweis IST der Eintrag. In einer Übersicht, die auf
%                  mehrere Stellen zeigt, steht daneben nichts, was die Nummer
%                  erklären würde; \ZSFref ergäbe dort eine Reihe blinder
%                  Sprungziele.
%
% Den Klartext liefert \nameref von selbst: Die Strukturmakros hinterlegen den
% Titel bereits bereinigt (\ZSFplainName, oben). Dieses Makro braucht deshalb
% keine eigene Liste stillzulegender Marker — die stünde sonst hier ein zweites
% Mal und veraltete beim nächsten neuen Marker.
\newcommand{\ZSFhubref}[1]{%
  \begingroup
  \ZSF@refComputeTargetColor{#1}%
  \hyperref[#1]{{\ZSFfontNote\ZSFfontRefLocator
    \ZSFInk{\ZSF@refTargetColor}{($\rightarrow$\,\ref*{#1}\,\nameref*{#1})}}}%
  \endgroup\allowbreak
}

% ------------------------------------------------------------
% Der Klartext-Name eines Ziels — die Grundlage von \nameref
% ------------------------------------------------------------
% Ein Titel darf Marker tragen: \SubsectionBar[sec:inverse]{Inverse
% \ZSFref{sec:hub-full-rank}}. Fuer die ANZEIGE des Balkens ist das richtig;
% als NAME des Labels ist es falsch — ein \ZSFref darin ergaebe einen Verweis
% im Verweis, ein \ZSFkeyword einen Registereintrag an der Fundstelle des
% Lesers statt an der des Ziels.
%
% Die Bereinigung steht deshalb an EINER Stelle und wird beim Setzen des
% Labels eingebacken (\ZSF@storeLabelName, aufgerufen aus den Strukturmakros).
% Jeder Konsument von \nameref bekommt sie automatisch — auch ein handge-
% schriebenes \nameref im Kapitel. Die Alternative waere, dass jede Lesestelle
% dieselbe Liste mitfuehrt; beim naechsten neuen Marker fehlt sie dort.
%
% \ZSFplainName traegt bewusst KEIN @: Der Name wandert ueber die .aux-Datei,
% und die wird ohne \makeatletter gelesen — gleiche Regel wie bei den
% Index-Makros \ZSFidxnum/\ZSFidxsee.
\NewDocumentCommand{\ZSF@plainKeyword}{s o m}{#3}
% Die Schlucker nehmen ihr Argument UND das Leerzeichen davor: Steht der Marker
% am Titelende (»Inverse \ZSFref{…}«), bliebe sonst ein Leerzeichen stehen.
\newcommand{\ZSF@eatOne}[1]{\unskip}
\newcommand{\ZSF@eatTwo}[2]{\unskip}
\NewDocumentCommand{\ZSFplainName}{+m}{%
  \begingroup
    \let\ZSFkeyword\ZSF@plainKeyword
    \let\ZSFlabel\@firstofone
    \let\ZSFhl\@firstofone
    \let\ZSFref\ZSF@eatOne
    \let\ZSFsectionref\ZSF@eatOne
    \let\ZSFScriptRef\ZSF@eatOne
    \let\ZSFTitleTag\ZSF@eatOne
    \let\ZSFindex\ZSF@eatOne
    \let\ZSFindexsee\ZSF@eatTwo
    #1%
  \endgroup
}
% Hinterlegt den Titel als Namen des folgenden \label. Ohne das liefert
% \nameref dokumentweit nichts — lautlos, weil kein Baustein es bisher las.
\newcommand{\ZSF@storeLabelName}[1]{\def\@currentlabelname{\ZSFplainName{#1}}}
\makeatother

\newcommand{\ZSFconclusion}[1]{%
  \ZSFInterlude{\ZSFspaceS}%
  \noindent$\Rightarrow$ #1%
  \ZSFInterlude{\ZSFspaceS}%
}

